% =====================================================================
%  chunk c09 —— 植田一石《ホモロジー的ミラー対称性》
%  第 19–24 节，原书 p033–p036
%  编号公式 34 条（p033: 7，p034: 8，p035: 9，p036: 10），逐页据 pages/pNNN.png 抄录
% =====================================================================
% [Tutor Note: 关于本 chunk 的公式条数。work/BRIEF.md 对 c09 记「50 条编号公式」，
%  但按页面图像逐页核对，p033–p036 上实际印出的编号公式共 34 条。BRIEF 的 50 应来自
%  work/tags_bypage.txt 的原始命中数（10+12+16+12=50），其中包含正文中的公式交叉引用
%  （如 p033 上的 (2.5) 出现 3 次、(21.1)/(21.2) 在 p034 正文中反复被引用、p035 上的
%  (21.6)、p036 上的 (23.1) 等）以及跨页重复计数，并非编号公式的条数。此处按图像
%  一一对应抄录 34 条，未增补任何编号，请 Lead 核定 BRIEF 的计数口径。]

\section{从同调镜像猜想走向 SYZ 猜想}

假设对 Calabi--Yau 流形的对 $(M, W)$，同调镜像对称 (2.5) 成立。于是对 $W$ 的每个点 $p$，存在 $M$ 的 Fukaya 范畴中对应于 $p$ 处的点层（skyscraper sheaf）$\cO_p$ 的对象。它一般而言是拉格朗日子流形的复形，但朴素地设想，假定它由 $M$ 的单个拉格朗日子流形实现。把这个拉格朗日子流形记作 $L_p$，则由范畴等价 (2.5) 得到
\begin{equation}
\mathrm{HF}^{*}(L_p) \cong \Ext^{*}(\cO_p, \cO_p) \tag{19.1}
\end{equation}
这里 $\mathrm{HF}^{*}(L_p)$ 是 $L_p$ 的 Floer 上同调。另一方面，$n$ 维复流形上点层的上同调总是同构于 $n$ 维环面 $T^n := (S^1)^n$ 的通常上同调；
\begin{equation}
\Ext^{*}(\cO_p, \cO_p) \cong H^{*}(T^n) \tag{19.2}
\end{equation}
在适当条件下，拉格朗日子流形的 Floer 上同调与通常上同调一致，因此若宽松地考虑，可以期望 $L_p$ 是环面。对相异的两点 $p, q \in W$，由范畴等价 (2.5) 得到
\begin{equation}
\mathrm{HF}^{*}(L_p, L_q) \cong \Ext^{*}(\cO_p, \cO_q) = 0 \tag{19.3}
\end{equation}
于是再次宽松地考虑，可以期望 $L_p$ 与 $L_q$ 不相交。进而，点层的集合 $\{\cO_p\}_{p \in W}$ 是张成类（spanning class，即对任意 $p \in W$，若 $\Hom(E, \cO_p) = 0$ 则 $E \cong 0$），因此自然期望，如此得到的这些 $L_p$ 铺满 $W$ 的整体，给出 $W$ 的拉格朗日环面纤维化。不过，Fukaya 范畴中的同构类在把拉格朗日子流形沿 Hamilton 同胚移动时不变，因此需要在拉格朗日子流形的 Hamilton 同痕类中取良好的代表元。类似的问题在复几何一侧也存在，构成模空间时的稳定性（stability）正对应于此。相干层的导出范畴中的稳定性，与 Calabi--Yau 流形的 flop 有很深的关联 [Bri02]。

\section{特殊拉格朗日子流形}

作为 Kähler 流形上的相干层，除点层之外还有全纯向量丛（holomorphic vector bundle）。由 Donaldson 与 Uhlenbeck--Yau 证明的小林--Hitchin 对应（Hitchin--Kobayashi correspondence），稳定的（stable）全纯向量丛与不可约 Hermite Yang--Mills 联络（irredubicle Hermitian Yang--Mills connection）成一一对应。在镜像一侧与之对应而受到期望的，正是特殊拉格朗日子流形。

% [Tutor Note: 上句中原文英文为 “irredubicle”，疑为 “irreducible” 之误，此处照录原文，请审校确认。]

我们曾把和乐群为 $\mathrm{SU}(n)$ 的子群的 $2n$ 维黎曼流形 $(M, g)$ 称为 Calabi--Yau 流形。$\C^n$ 上的 $(1, 1)$ 形式
\begin{equation}
\omega_0 := \frac{\sqrt{-1}}{2}\left(dz_1 \wedge d\bar z_1 + \cdots + dz_n \wedge d\bar z_n\right) \tag{20.1}
\end{equation}
与 $(n, 0)$ 形式
\begin{equation}
\Omega_0 := dz_1 \wedge \cdots \wedge dz_n \tag{20.2}
\end{equation}
在 $\mathrm{SU}(n)$ 的作用下不变，因此从某点的切空间出发、通过平行移动加以延拓，便可得到 Calabi--Yau 流形上的微分形式 $\omega$ 与 $\Omega$，使其共变微分为零。它们分别给出 $M$ 的 Kähler 形式与全纯体积形式（holomorphic volume form）。$\omega_0$ 与 $\Omega_0$ 满足关系式
\begin{equation}
\frac{1}{n!}\omega_0^{n} = (-1)^{n(n-1)/2}\left(\frac{\sqrt{-1}}{2}\right)^{n}\Omega_0 \wedge \bar\Omega_0 \tag{20.3}
\end{equation}
因此 $\omega$ 与 $\Omega$ 也同样满足
\begin{equation}
\frac{1}{n!}\omega^{n} = (-1)^{n(n-1)/2}\left(\frac{\sqrt{-1}}{2}\right)^{n}\Omega \wedge \bar\Omega \tag{20.4}
\end{equation}
$M$ 的 $n$ 维子流形 $L$ 满足
\begin{equation}
\omega|_{L} = 0 \tag{20.5}
\end{equation}
时，称 $L$ 为拉格朗日子流形（Lagrangian submanifold）。若存在某个实数 $\theta$，使拉格朗日子流形 $L$ 满足
\begin{equation}
\operatorname{Im}\left(\ee^{\sqrt{-1}\theta}\Omega\right)\Big|_{L} = 0 \tag{20.6}
\end{equation}
则称 $L$ 为相位（phase）为 $\theta$ 的特殊拉格朗日子流形（special Lagrangian submanifold）。把 $\Omega$ 乘以常数即总可使相位为 $0$，因此只要处理单个特殊拉格朗日子流形，相位就不是本质的。

\section{极小子流形}

设 $(M, g_M)$ 为黎曼流形，$L$ 为其子流形。$L$ 的切丛关于 $g_M$ 的正交补空间称为 $L$ 的法丛（normal bundle），记作 $N_{L/M}$；
\begin{equation}
TM_{L} := TL \perp N_{L/M}. \tag{21.1}
\end{equation}
对 $L$ 上的向量场 $v, w \in \Gamma(TL)$，任意取 $M$ 上的向量场 $\tilde v, \tilde w \in \Gamma(TM)$ 使 $\tilde v|_{L} = v$ 且 $\tilde w|_{L} = w$，并置
\begin{equation}
B(v, w) := \pi_{N_{L/M}}\left(\nabla_{\tilde v}\tilde w\right) \tag{21.2}
\end{equation}
其中，
\begin{equation}
\pi_{N_{L/M}} : TM_{L} \to N_{L/M} \tag{21.3}
\end{equation}
是关于正交分解 (21.1) 的投影，$\nabla$ 是 $TM$ 关于 $g_M$ 的 Levi-Civita 联络。(21.2) 的右端显然不依赖于 $\tilde v$ 的取法。又，由 $\tilde v|_{L}, \tilde w|_{L} \in \Gamma(TM|_{L})$ 是 $\Gamma(TL)$ 的元素可知 $[\tilde v, \tilde w]|_{L} \in \Gamma(TM|_{L})$ 也是 $\Gamma(TL)$ 的元素，再由 Levi-Civita 联络的零挠率条件 $\nabla_{\tilde v}\tilde w - \nabla_{\tilde w}\tilde v = [\tilde v, \tilde w]$，可知 (21.2) 的右端关于 $\tilde v$ 与 $\tilde w$ 的交换是对称的。因此 (21.2) 的右端也不依赖于 $\tilde w$ 的取法，它确定 $\Gamma(\mathrm{Sym}^{2}\,T^{*}L \otimes N_{L/M})$ 的一个元素，称之为 $L$ 的第二基本形式（second fundamental form）。关于 $L$ 的度量 $g_L := g|_{L}$ 的第二基本形式 $B$ 的迹（trace），称为 $L$ 的平均曲率（mean curvature）；
\begin{equation}
H := \operatorname{tr}_{g|_{L}} B \in C^{\infty}(L, N_{L_M}). \tag{21.4}
\end{equation}
关于体积泛函
\begin{equation}
\operatorname{Vol}(L) := \int_{L}\operatorname{vol}L \tag{21.5}
\end{equation}
的 Euler--Lagrange 方程，已知给出二阶非线性偏微分方程
\begin{equation}
H = 0 \tag{21.6}
\end{equation}

% [Tutor Note: (21.4) 中 $C^\infty$ 的第二个变量原文印作 $N_{L_M}$（M 为 L 的下标），
%  与 (21.2)(21.3) 的 $N_{L/M}$ 不一致，疑为排版时漏掉斜线；此处照录图像，请审校确认。]

\section{校准几何}

黎曼流形 $(M, g_M)$ 上的 $k$ 次微分形式若满足以下两个条件，则称为 $k$ 次校准形式（calibration）：
\begin{enumerate}[label=\arabic*.]
\item $\varphi$ 是闭形式。
\item 对任意点 $p \in M$ 与任意有向的 $k$ 维子空间 $W \subset T_pM$，不等式
\begin{equation}
\varphi\big|_{\wedge^{k}W} \le \operatorname{vol}W \tag{22.1}
\end{equation}
成立。
\end{enumerate}
若 $M$ 的有向 $k$ 维子流形 $N$ 对任意 $p \in N$ 都满足不等式 (22.1) 中的等号
\begin{equation}
\varphi\big|_{\wedge^{k}T_pN} = \operatorname{vol}_{T_pN} \tag{22.2}
\end{equation}
则称 $N$ 为校准子流形（calibrated submanifold）。由校准形式的定义立即可得不等式
\begin{equation}
\operatorname{Vol}(N, g|_{N}) = \int_{N}\operatorname{vol}_{N} \le \int_{N}\varphi \tag{22.3}
\end{equation}
注意，(22.3) 的右端只依赖于 $N$ 的同调类 $[N] \in H_k(M, \Z)$。这是关于 Yang--Mills--Higgs 泛函的 Bogomolnyi--Prasad--Sommerfield 不等式的类似物。由不等式 (22.3) 立即推出，校准子流形在其所属的同调类中具有最小体积。特别地，校准子流形是极小子流形。极小子流形的方程 (21.6) 与校准子流形的方程 (22.2) 之间的关系，与 Yang--Mills 方程 $D^{*}F = 0$ 和 ASD 方程 $*F = -F$ 之间的关系十分相似。前者是二阶偏微分方程，给出泛函的驻点；后者是一阶偏微分方程，给出泛函的最小值。

校准形式与校准子流形的概念是由 [HL82] 引入的。校准形式最基本的例子是 Kähler 流形中 Kähler 形式的幂 $\omega^{k}$；此时成为校准子流形的充要条件是：它是复子流形。对 Calabi--Yau 流形 $(M, g, J, \omega, \Omega)$，$\operatorname{Re}\Omega$ 是校准形式，此时的校准子流形是特殊拉格朗日子流形。其他校准子流形的例子，还有 $G_2$ 流形中的结合子流形（associative 3-fold）、余结合子流形（coassociative 4-fold），以及 $\mathrm{Spin}(7)$ 流形中的 Cayley 子流形（Caley 4-fold）等。

% [Tutor Note: 原文 “Caley 4-fold” 疑为 “Cayley 4-fold” 之误，此处照录原文，请审校确认。]

\section{McLean 定理}

设 $(M, g_M, J, \omega, \Omega)$ 为 Calabi--Yau 流形，$L$ 为其特殊拉格朗日子流形。关于 $L$ 的黎曼度量 $g_L := (g_M)|_{L}$ 的 Hodge 星算子记作 $*_{L} : \Omega^{k}(L) \to \Omega^{n-k}(L)$。设 $\nu \in \Gamma(N_{L/X})$ 为 $L$ 的法向量场，则
\begin{equation}
(\iota_{\nu}\operatorname{Im}\Omega)\big|_{L} = -*_{L}\left((\iota_{\nu}\omega)\big|_{L}\right) \tag{23.1}
\end{equation}
成立。事实上，(23.1) 不含微分，因此逐点取切空间的适当基，即可归结为 $M = \C^{n}$ 且 $L = \R^{n}$、并且
\begin{align}
\Omega &= dz_1 \wedge \cdots \wedge dz_n = (dx_1 + \sqrt{-1}dy_1) \wedge \cdots \wedge (dx_n + \sqrt{-1}dy_n), \tag{23.2}\\
\omega &= dx_1 \wedge dy_1 + \cdots + dx_n \wedge dy_n, \tag{23.3}\\
\nu &= \frac{\partial}{\partial y_1} \tag{23.4}
\end{align}
的情形；此时容易看出 (23.1) 的两端都是 $dx_2 \wedge \cdots \wedge dx_n$。

$L$ 沿法向量场 $\nu$ 的微小变形成为特殊拉格朗日子流形的充要条件是
\begin{equation}
\cL_{\nu}\operatorname{Im}\Omega = \cL_{\nu}\omega = 0 \tag{23.5}
\end{equation}
利用 Cartan 的同伦公式
\begin{equation}
\cL_{\nu} = d \circ \iota_{\nu} + \iota_{\nu}d \tag{23.6}
\end{equation}
以及 $\Omega$ 和 $\omega$ 是闭形式这一点，此条件可改写为
\begin{equation}
d\iota_{\nu}\omega = d\iota_{\nu}\operatorname{Im}\Omega = 0 \tag{23.7}
\end{equation}
由 (23.1)，此条件又等价于
\begin{equation}
d\iota_{\nu}\omega = d(*\iota_{\nu}\omega) = 0, \tag{23.8}
\end{equation}
即 $\iota_{\nu}\omega$ 是调和 $1$ 形式。通过使用 Banach 流形上的隐函数定理，McLean 证明了下面的定理：

\begin{theorem}[{[McL98, Theorem 3.6]}]
特殊拉格朗日子流形的模空间是光滑的，且其在 $L$ 处的切空间自然同构于 $L$ 上的调和形式空间。
\end{theorem}

特别地，特殊拉格朗日子流形的模空间的维数等于 $L$ 的第一 Betti 数。对 $\nu_1, \nu_2 \in \Gamma(N_{L/M})$，置
\begin{equation}
g(\nu_1, \nu_2) = \int_{L_b}(\iota_{\nu_1}\omega) \wedge *(\iota_{\nu_2}\omega) \tag{23.9}
\end{equation}
则它给出特殊拉格朗日子流形的模空间的度量，称之为 McLean 度量（McLean metric）。

% [Tutor Note: (23.1) 上方的 ν ∈ Γ(N_{L/X}) 中原文用 X，与后文 Γ(N_{L/M}) 不一致，
%  疑为 M 之误；此处照录图像。(23.9) 的积分区域原文印作 L_b（b 为下标），
%  疑为 L 之误或排印残留；此处照录图像。]

\section{Hessian 结构}

仿射流形 $M$ 上的黎曼度量 $g$ 称为 Hessian 度量（Hessian metric），是指：对每点 $p \in M$，存在 $p$ 的邻域 $U$ 的仿射坐标 $\{y_i\}_{i=1}^{n}$ 与 $U$ 上的函数 $K$，使得
\begin{equation}
g_{ij} = \frac{\partial^{2}K}{\partial y_i\partial y_j} \tag{24.1}
\end{equation}
成立。用确定 $M$ 的仿射结构的联络 $D$，(24.1) 可写成不依赖坐标的形式
\begin{equation}
g = DdK \tag{24.2}
\end{equation}
注意，一般地 $D$ 与 $g$ 的 Levi-Civita 联络 $\nabla$ 并不一致。由于度量 $g$ 正定，$K$ 成为凸函数，并且除加仿射函数的自由度外被唯一确定。

用仿射坐标把 $U$ 与 $\R^{n}$ 的开集等同，并定义映射 $\check y = (\check y_1, \dots, \check y_n) : U \to \R^{n}$ 为
\begin{equation}
\check y_i = \frac{\partial K}{\partial y_i} \tag{24.3}
\end{equation}
则由 $K$ 的凸性，$\check y$ 是单射，并给出 $U$ 的坐标。进而，定义 $U$ 上的函数 $\check K$ 为
\begin{equation}
\check K := \sum_{i=1}^{n}\check y_i y_i - K \tag{24.4}
\end{equation}
则
\begin{align}
\frac{\partial\check K}{\partial\check y_i} &= y_i + \sum_{j=1}^{n}\check y_j\frac{\partial y_j}{\partial\check y_i} - \frac{\partial K}{\partial\check y_i} \tag{24.5}\\
&= y_i + \sum_{j=1}^{n}\frac{\partial K}{\partial y_j}\frac{\partial y_j}{\partial\check y_i} - \frac{\partial K}{\partial\check y_i} \tag{24.6}\\
&= y_i \tag{24.7}
\end{align}

% 第 24 节正文续 p037，由 c10 接续。

% CHUNK c09 DONE pages p033–p036
